\documentclass[a4paper,10pt]{report}
\usepackage{graphicx}
\usepackage{hyperref}

\newcommand{\gvmtic}{\textsc{gvmtic}}
\newcommand{\gvmtc}{\textsc{gvmtc}}
\newcommand{\glink}{\textsc{gvmtlink}}
\newcommand{\gvmtas}{\textsc{gvmtas}}
\newcommand{\gvmtcc}{\textsc{gvmtcc}}
\newcommand{\gvmtxc}{\textsc{gvmtxc}}

% Title Page
\title{The Glasgow Virtual Machine Toolkit Manual}
\author{Mark Shannon}
\date{ } 

\begin{document}
\maketitle 

\begin{abstract}
The GVMT is a toolkit for building virtual machines. The GVMT does not dictate the overall design of the virtual machine, but it does eliminate a lot of implementation details. 

The developer needs to develop an interpreter and supoorting code, as normal. The toolkit handles garbage collection transparently
and can automatically create a just-in-time compiler from the interpreter source code.

This manual covers the components of the Glasgow Virtual Machine Toolkit, hereafter called the GVMT, and their use. 

Convertion of source-code to bytecodes is not handled by the GVMT, that is up to the developer.

\tableofcontents

\end{abstract}


\chapter{Introduction}

The Glasgow Virtual Machine Toolkit, GVMT, is a toolkit for constructing high-performance virtual machines. 

The GVMT can be divided into three parts: \begin{itemize}
\item Tools for the creation of interpreters, compilers and other virtual machine components.
\item Library code. Primarily for Garbage Collection.
\item An interface between the user-defined parts of the virtual machine and the toolkit library.
\end{itemize}

All the components operate on a common abstract machine model. It is not necessary to understand this model to use the GVMT, but a knowledge of the abstract machine is useful to get the best out of the GVMT. The abstract machine, which is described in detail in Chapter \ref{chap:abstract}, is a stack-based machine supporting concurrency and garbage-collection. It has an extensible instruction set. The instruction set is abstract in that it cannot be directed executed, but is used to defined the semantics of bytecodes and other code. 

The tools provided by the GVMT can be split into front-end tools which compile source code to the abstract machine model, and back-end tools which convert from the abstract machine code to real machine code.


\section{Example virtual machine}

In order to help understand the GVMT, a simple Scheme implementation is included as an example, it can be found in the \verb|/example| directory. 
This example illustrates all the important aspects of the GVMT, while not being overly complex.

All scheme VMs are implemented by what is known as read-eval-print loop. Its reads the input, evaluates the input and then prints the result, repeating until the interpreter exits. 
The `read' part consists of reading text from a file or terminal and parsing that text to form a syntax tree.

The file \verb|/example/parser.c| contains the parser, it is a standard C(89) source file which can be compiled with the GVMT C compiler, \gvmtc (see Section \ref{sect:gvmtc})

In some scheme implementations the `eval' part is implemented by directly evaluating the syntax tree. This is not the case for the GVMT version, as all GVMT-based VM use bytecode interpreters. The GVMT scheme implementation first compiles the syntax tree to bytecode, then evaluates the resulting bytecode.
The syntax-tree to bytecode compiler code can be found in \verb|/example/compiler.c|.


The bytecode interpreter is defined in \verb|/example/interpreter.vmc|

Although converting to bytecode is extra work it brings two benefits. The first is that the resulting bytecode can be optimised more easily than optimising the syntax tree directly. The second is that the GVMT can produce a bytecode to machine-code compiler automatically.

The Scheme optimiser is a series of passes, which use special interpreters created using \gvmtxc\ (see Section \ref{sect:gvmtxc}).
These passes are defined in the files of the form \verb|/example/XXX.vmc| (except \verb|/example/interpreter.vmc|).

The bytecode to machine-code, or JIT,\footnote{JIT is short for Just-In-Time and is the standard term for runtime compilers} compiler is generated by the \gvmtcc\ tool from the output of \gvmtic.

\section{The Tools}

The GVMT provides a number of tools, which are described in more detail in Chapter \ref{chap:tools}. These can be loosely grouped into front-end tools, and back-end tools. Front-end tools are responsible for compiling source code to the GVMT abstract machine code. Back-ends convert this abstract machine code into an executable virtual machine. 

\subsection{The front-end compilers, \gvmtc\ and \gvmtic}
\gvmtc{} is a C compiler and compiles C files to GSC files.
\gvmtic{} Takes an interpreter specification file and converts into a GSC file.
The interpreter specification is in the form of a list of bytecodes; each of which consists of a name, a stack effect comment and a C code (or GSC) body.

Both tools use a modified version of lcc (lcc-gvmt) to do the compilation. See section \ref{sect:lcc} for more details.

\begin{figure}[htb]
 \centering
 \caption{Building a Virtual Machine using GVMT\label{fig:build}}
 \includegraphics[scale=0.4, bb=36 36 900 700]{build.eps}
\end{figure}

\subsection{The code-generators, gvmtas and gvmtcc}
The code-generators take GSC as input and create object files or executables.

\gvmtas{} compiles the GSC into the GVMT object format, GSO. It works by converting GSC back into C in such a way as to respect the semantics of the GVMT execution model, on a section-by-section basis. \gvmtas{} is also responsible for wrapping the bytecodes in an interpreter loop. 

\gvmtcc{} takes a GSC file containing a bytecode section and produces a bytecode to machine-code. Currently \gvmtcc\ generates llvm-IR\footnote{See llvm.org for more details} and uses LLVM to generate machine-code at runtime (as a JIT compiler).

\subsection{Other tools}

\gvmtxc\ takes a GSC file containing a bytecode section representing the main interpreter, plus a secondary file describing a secondary interpreter. It produces an interpreter (as a GSC file) which uses the bytecodes defined for the main interpreter, but with the semantics defined in the secondary file. \gvmtxc{} can be used for creating bytecode verifiers and disassemblers, as well as optimisers and other abstract interpreters.

The bytecode-transformation engines produced by \gvmtxc{} and the compilers produced by \gvmtcc{} can be used independently or combined to form a specialised optimising compiler.

\glink\ links together all the GSO files produced by \gvmtas\ into a single object ready to be linked by the system linker. See Figure \ref{fig:build}.

\section{GVMT input files}
GVMT takes a number of input files:
\begin{itemize}
\item Interpreter description file(s).
\item Function definition file(s).
\end{itemize}




\chapter{The tools\label{chap:tools}}



GVMT is, as the name suggests, supplies a number of tools. Six of these form the core of the toolkit. 

As mentioned in the previous chapter, the GVMT tools can be grouped into front-end and back-end tools. The front-end tools convert source code into code for the GVMT Abstract Machine. Back-end tools convert the abstract representation into real manchine code.

\section {Core tools}
The six core tools are:
\begin{itemize}
 \item Front end tools
\begin{itemize}
\item \gvmtic
\item \gvmtc
\item \gvmtxc
\end{itemize}
 \item Back end tools
\begin{itemize}
\item \gvmtas
\item \gvmtcc
\item \glink
\end{itemize}
\end{itemize}

\gvmtic{} and \gvmtc{} form the front-end and are used to translate the C-based interpreter description and standard C files, respectively, into GSC format.

\gvmtas{} converts these GSC files to object files and \glink{} links them together.
\gvmtcc{} takes the GSC file produced by \gvmtic{} and produces a compiler from it.

\subsection{Interpreter description file}
An interpreter description file consists of bytecode definitions plus the local variables shared by all bytecodes.

Each bytecode can be defined by its stack effect plus a snippet of C code or as a sequence of other bytecodes and predefined GVMT abstract stack machine codes.

Usually there is one primary interpreter description file that defines the bytecode format. Secondary interpreters (see Sect. \ref{sect:gvmtxc} will be guaranteed to use the same bytecode format. However it is possible to have several independent primary interpreters. For example a virtual machine might use one bytecode format for its interpreter and compiler and an entirely different format for its regular-expression engine. Each primary interpreter may have secondary interpreters.

Interpreters are fully re-entrant and thread-safe. The interpreter description file  defines not only the behaviour of the interpreter, but the behaviour of the GVMT-generated compiler as well.

\subsubsection{Format}

The interpreter description file consists of a locals block and any number of bytecode declarations in any order. 
\subsubsection{Interpreter-local variables}
Interpreter-locals variables are declared in a locals section:
\begin{verbatim}
locals {
        type name;
\end{verbatim}
\nopagebreak
        \vdots
\begin{verbatim}
}
\end{verbatim}

Interpreter-local variables are visible within all bytecode definitions, and can be accessed indirectly via the interpreter frame on the control stack. Interpreter-local variables (and the interpreter frame) have the same lifetime as the invocation of the interpreter.

\subsubsection{Bytecode definitions}
Bytecodes can be defined either in C or directly in GSC. Both types of definition can be used in the same interpreter.

C-bytecode definitions are in the form

\begin{verbatim}
name (`[' qualifier* `]')?  `(' stack_comment `)' [ `=' <int> ] `{'
    Arbitrary C code
`}'
\end{verbatim}

Legal qualifiers are:
\begin{description}
\item [private] This bytecode is not visible at the interpreter level.
\item [protected] This bytecode is acceptable to the interpreter, but not to any verification code. It is only to be used internally as an optimisation. The meaning of `protected' is largely conventional. GVMT will ignore it.
\item [nocomp] This bytecode will not appear in bytecode passed to the compiler. If this bytecode is seen by the compiler it will abort.
Useful for reducing the size of the compiler.
\item [componly] This bytecode will only be passed to the compiler. It will cause the interpreter to abort. Useful for reducing the size of the interpreter.
\end{description}

The \verb|stack_comment| is of the form:
 \begin{verbatim}
type id (`,' type, id)* -- type id (`,' type, id)*)
\end{verbatim}

Where \verb|type| is any legal C type and \verb|id| is any legal C identifier, or a legal C identifier prefixed with \verb|#|. Variables prefixed with \verb|#| are fetched from the instruction stream, not popped from the stack. When referred to in the C code, the \verb|#| prefix is dropped.

The integer at the end of the header line is the actual numerical value of this bytecode. If not supplied, GVMT will allocate bytecode numbers starting from 1.

Compound instructions are in the form:
\begin{verbatim}
name (`[' qualifier* `]')?  '(' stack comment ')' [ `=' <int> ] `:'
   instruction*
`;'
\end{verbatim}

Where \verb|instruction| is a legal GSC or user-defined instruction.

To allow user defined instructions to do anything useful, a large number of builtin instructions are provided. All are private, that is they must be used inside a compound instruction to be visible at the interpreter level. 

Instruction stream manipulation operators are provided:
\verb|#@| Removes a byte from the instructino stream and pushes it to the data stack.
\verb|#N| where N is a one-byte integers pushes N into the front of the instruction stream.
\verb|#[N]| Pushes a copy of the N\textsuperscript{th} item, starting at zero, in the instruction stream into the front of the instruction stream.
\verb|#+| or \verb|#-| Adds or subtracts, respectively the first two values from the instruction stream  and pushes the result back to the front of the instruction stream.

Appendix \ref{app:inst} contains a full list of all abstract machine instructions 

There are three bytecode names which are treated specially by the GVMT, these are:
\begin{itemize}
\item \_\_preamble 
\item \_\_postamble
\item \_\_trace
\end{itemize}
\verb|__preamble| is executed when the interpreter is called, before any bytecodes are executed.

\verb|__postamble| is executed when the interpreter reaches the end of the bytecodes. Interpreters with a \verb|__postamble| instruction expect a second parameter, the end of the bytecode sequence. This is useful for `interpreters' which analyse code rather than executing it.

\verb|__trace| is executed before the main bytecode definition for every bytecode executed.


\subsubsection*{Examples}


TO DO -- Put a couple of examples here. 1 C, 1 vmc.



\subsection{The interpreter generator \gvmtic{}\label{sect:gvmtic}}
\gvmtic{} comiles an interpreter description file, as described above, into a GSC file.

\gvmtic{} takes the following options:

\input{gvmtic_options.tex}

\subsection{The C compiler \gvmtc{}\label{sect:gvmtc}}
\gvmtc{} takes a standard C89 file as its input and outputs a GSC file.

\gvmtc{} takes the following options:

\input{gvmtc_options.tex}

\subsection{The GSC assembler \gvmtas{}\label{sect:gvmtas}}
\gvmtas{} takes a GSC file and assembles it to a GVMT object file (.gso is the expected file extension)

\gvmtas{} takes the following options:

\input{gvmtas_options.tex}

\subsection{The compiler generator \gvmtcc{}\label{sect:gvmtcc}}
\gvmtcc{} Takes a GSC file produced by \gvmtic{} and produces a compiler. The generated compiler is in C++ and relies on LLVM (http://llvm.org/).
It will need to be compiled with the system C++ compiler.

It may be possible to generate compilers using other code-generator back-ends in the future, but there is currently no plan to do so.

\gvmtcc{} takes the following options:

\input{gvmtcc_options.tex}

\subsection{The linker \glink{}\label{sect:glink}}
\glink{} takes any number of GVMT object (.gso) files and links them to form a single native object file.
This object file can be linked with the GC object file and compiler object file (if required)
using the system linker to form an executable.

\glink{} takes the following options:

\input{gvmtlink_options.tex}

\subsection{lcc-gvmt\label{sect:lcc}}
Both \gvmtic{} and \gvmtc{} use a modified version of lcc as their C-to-GSC compiler. This compiler uses the standard lcc front-end and a custom back-end. The back-end operates in four passes. The first pass labels the IR forests with the C type for that expression, as much as can be derived without the explicit casts present in the source code. The second pass is a bottom pass which labels the trees with the set 
of possible GSC types it may take. The third pass is a top-down pass which attempts to find the exact GSC type of each node. The fourth pass emits the GSC.

\section{Other tools}

\subsection{The bytecode-processor generator \gvmtxc\label{sect:gvmtxc}}
\gvmtxc\ Takes a (partial) interpreter description file and a master GSC file (produced by \gvmtic) and produces a GSC file representing an interpreter. This interpreter is guaranteed to accept the same bytecode format as the master interpreter.
For a partial interpreter description file, the generate interpreter skips over the omitted bytecodes.

\gvmtxc{} takes the following options:

\input{gvmtxc_options.tex}


\subsection{The object layout tool\label{sect:objects}}

The object layout tool is an optional tool to help manage heap object layout. It can create C structs, shape vectors and marshalling to GSC files for objects all from a single specification.
This tool is designed as an extensible Python script rather than a stand-alone tool, although it can be used as such.
When used as a standalone tool it has takes the following options:

\input{objects_options.tex}


\chapter{GVMT Execution Model\label{chap:abstract}}


The GVMT assists in building a virtual machine(VM) for a high-level language(HLL). 
Like any virtual machine, the user-defined VM must execute on a lower-level machine (usually a hardware machine). The GVMT defines an abstract stack-based machine (The GVMT Abstract Machine) that has well defined semantics, but is abstract in the sense that it cannot execute any programs. It is designed to both assist in implementing stack-based virtual machine and allow translation to efficient machine-code.

\section{Components}
The GVMT Abstract Machine consists of two parts:

\begin{enumerate}
\item A set of executing threads
\item A shared heap
\end{enumerate}

Each thread of execution consists of five components:
\begin{enumerate}
\item Temporary registers
\item Control stack
\item Data stack
\item Thread-local variables
\item Exception handling stack
\end{enumerate}

When the GVMT Abstract Machine is started it consists of one thread only.

\section{Threads}
Each thread is independent, the number of concurrently executing threads is limited only by the hardware. Threads share the heap, but have their own stacks, thread-local variables and temporary registers.

\subsection{Execution Model}
The execution model executes abstract HLL virtual machine instructions. Each HLL VM instruction is defined in terms of GSC instructions.
GSC instructions can be grouped into three categories:
\begin{description}
\item[Data] Instructions that take values on the Data stack and place results there, such as addition.
\item[Flow control] Instructions that alter the flow of control, both within HLL-VM instructions and between HLL-VM instructions.
\item[Stack manipulation] Instructions that explicitly manipulate or describe the status of the data and control stacks.
\end{description}

Appendix \ref{app:inst} contains a list of all the GSC instructions.

\subsection{GC Safe points}
The GC safe instruction declares that the executing thread can be interupted for garbage collection. In order for this to be safe, the following restrictions must be observed:
\begin{enumerate}
\item No non-reference values are on the execution stack.
\item All heap objects reachable from this thread are in a scannable state.\footnote{Usually this means that the header(class) of any newly created object will have been initialised.}
\end{enumerate}
The front-end tools \gvmtc{} and \gvmtic{} automatically ensure that condition 1 is true for all C code.
The toolkit ensures that all virtual stack items are reachable by the garbage-collector.
Condition 2 can be ensure simply by initialising all objects immediately after allocation.

\subsection{The control stack}
The control stack is an opaque structure used to handle procedure calls and returns, exception handling and storing of temporary variable, across calls. Interpreter frames are also allocated on the control stack. Memory can be allocated on the call-stack, but subsequent allocation are not guaranteed to be contiguous. Parameters are also passed to native code using the control stack.
The control stack will usually map directly onto the native C stack.

\subsection{Data stack}

        \begin{center}
        \includegraphics[scale = 0.8]{Stack.eps}
% Diagram1.eps: 300dpi, width=3.29cm, height=4.97cm, bb=0 0 388 587
        \end{center}


The data stack is where all arithmetic operations takes place. The top part of the data stack may be virtualised, that is stored in registers or other discrete values. The data stack is continuous across calls, and is used for passing parameters. Parameters are pushed to the stack by the caller and popped by the calle. See the GSC instruction \verb|STACK| for more details on how to access the in-memory stack directly.

\subsection{Thread-local variables}
Each thread has thread-locals variables, the number and type of these variables are determined by the developer.

\section{The heap}
The GVMT Abstract Machine heap is a garbage collected heap. The toolkit user needs to provide a function \verb|gvmt_shape| which supplies the garbage collector with information for scanning of objects. See section \ref{sect:user-shape}.
Static roots of the heap can either be defined when the VM is built or added at runtime, see section\ref{sect:user-roots}.

\subsection{Shape\label{sect:shape}}
During tracing the garbage collector needs to know both the extent of an object and which fields are references to other objects, and which are not. The developer communicates this information to the toolkit via a `shape' vector (array).
This shape is a zero-terminated vector of integers. Each integers represents either N successive references or N successive words of non-references. References are represented by positive numbers, non-references by negative numbers (zero is the terminator).
For example the shape [ 2, -2, 1, 0 ] represents an object of total extent 5 words, the first 2 and last of which are references.
The following C declaration (taken from the example Scheme VM):
\begin{verbatim}
GVMT_OBJECT(cons) {   
    struct type *type;   // Opaque (non-GC) pointer
    GVMT_Object car;     // Reference to heap object
    GVMT_Object cdr;     // Reference to heap object
};
\end{verbatim}
defines \verb|cons| objects. These \verb|cons| objects would have a shape of [ -1, 2, 0 ].

\subsection{Garbage collection}
No garbage collection algorithm is specified for the GVMT. However, since the GVMT can support both read and write barriers as well as the declaration of GC safe points, a wide range of collectors should be feasible.
Currently both generational and non-generational collectors are available, although only `stop-the-world' collectors have been implemented. The GVMT garbage collectors also support tagging, which allows non-references to be stored in reference slots by setting one or more `tag' bits.
Read and write barriers are implemented via the RLOAD\_R and RSTORE\_R instructions respectively. 

\subsection{Exception handling}
The GVMT exception handling model uses an explicit stack, rather than tables. These means that using exceptions has a runtime cost even when the exceptions are not used as a means of flow control transfer, but that exception handling itself is much faster than for table-based approaches. The GVMT exception handling mechanism is similar to the C setjump-longjump mechanism.
The six GSC instructions involved are (See section \ref{sect:user-except} for the matching intrinsics)
\begin{itemize}
\item [PUSH_CURRENT_SATE] Pushes the current execution state onto the state stack. This state-object holds the following information: The current execution point(immediately after the PROTECT instruction), the stack depth, and the current control stack frame. Finally it pushes NULL onto the data stack. 
\item [DISCARD_STATE] Pops the state-object from the state stack and destroys it.
\item [PUSH_STATE] Pushes a previously popped state-object to the state stack.
\item [POP_STATE] Pops state-object from the state stack, leaving it on the data stack. State-objects are \emph{not} garbage collected. 
Use  DISCARD_STATE to dispose of state-objects.
\item [RAISE] Pops the value from the TOS and stores it. The state-object on top of the state stack is examined and the thread execution state (data and control stack pointers) restored to the values held in the handler. The stored value is pushed onto the data stack. Finally, execution resumes from the location stored in the state-object.
\item [TRANSFER] Pops the value from the TOS and stores it. The state-object on top of the state stack is examined and the thread execution state, except the data-stack (control stack pointer only) restored to the value held in the handler. The stored value is pushed onto the data stack. Finally, execution resumes from the location stored in the state-object.
\end{itemize}
It is an error (resulting in a probable crash) to leave an exception handler on the exception stack after the procedure in which it was pushed has returned.
There is no concept of catching only some sorts of exceptions in the GVMT. All exceptions are caught; exceptions that should be handled elsewhere must be explicitly reraised. 

Exceptions work as expected in interpreted (and compiled) code. The interpreter instruction pointer is stored along with the machine instruction pointer.

\section{Concurrency\label{sect:conc}}

The GVMT is concurrency friendly, all generated code and library code is thread-safe. The GVMT does not provide a thread library, so the developer will have to use the underlying operating system library.
Concurrency is at the thread level, the GVMT is unaware of processes.

In order that the toolkit can operate correctly, it must be informed of the creation and destruction of threads.
The toolkit provides functions for this purpose, see Section \ref{sect:user-threads} for more details.
The developer should also be aware that heap allocation may block, waiting for a collection, so locks should not be held across an allocation.

The GVMT also provides fast locks for synchronisation in the case where contention is expected to be rare. This functionality is provided by the  LOCK and UNLOCK abstract machine instructions. 

The intrinsic functions \verb|gvmt_lock(gvmt_lock_t *lock)| and \verb|gvmt_unlock(gvmt_lock_t *lock)| are translated to the LOCK and UNLOCK instructions, respectively.

\emph{Warning:} GVMT locks are not automatically unlocked by the RAISE instruction. The developer must ensure that either a RAISE cannot occur between a LOCK and UNLOCK, or that the LOCK UNLOCK pair are PROTECTed. A safe LOCK/UNLOCK sequence is:
\begin{verbatim}
gvmt_lock_t lock = GVMT_LOCK_INITIALIZER;
GVMT_Object ex;
gvmt_lock(&lock)
GVMT_BEGIN_TRY(ex)
\end{verbatim}
\nopagebreak[4]
        \vdots
\nopagebreak[4]
\begin{verbatim}
gvmt_unlock(&lock)
GVMT_CATCH
gvmt_unlock(&lock)
gvmt_raise(ex)
GVMT_END_TRY
\end{verbatim}

\subsection{Memory model}
As all threads share the same heap, they may access heap objects concurrently. The following guarantees are made about the state of the heap and roots seen from one thread when modified by another thread:

All writes to pointers(P), references(R) and integers(I) not larger than a pointer are atomic.


The order in which changes to memory appear to one thread are the same as for another. Ie, thread 1 writes to A, then writes to B. Thread 2 will not see a change in B before it sees a change in A.


In between synchronisation points, the delay between a write being made by one thread and a write being seen by another thread is indefinite.

\section{GVMT Stack Code}
GVMT Stack Code (GSC) is the instruction set of the GVMT Abstract Machine.
GSC should be viewed as a sort of abstract assembly language. It is possible to define a virtual machine entirely in C, as the front-end tools, \gvmtic{} and \gvmtc, compile C into GSC.
It is useful, however to have some awareness of GVMT Stack Code when using the GVMT.
To get a better idea of how GSC is used, have a look at the output of \gvmtc and \gvmtic in the \verb|/example| directory.

Appendix \ref{app:inst} contains a full list of GSC instructions.


\chapter{Building a VM using the Toolkit}

\section{Before you start}
Before starting a virtual machine, you will need to generate some bytecode programs to execute. If you are targetting a pre-existing VM format, this is straightforward. If you are creating a new language you will need a parser. \gvmtic{} can produce a C header file containing definitions for all the opcodes: the C macro for the opcode for instruction `xxx' is \verb|GVMT_OPCODE(xxx)|.

\section{Defining the components}
For GVMT to build a VM, you will need to define:
\begin{itemize}
\item An interpreter. See section \ref{sect:gvmtic}
\item Any other support code that your interpreter calls. See section \ref{sect:gvmtc}
\end{itemize}
You will also need to define a small number of functions required by the GVMT. The header file for these can be found in /include/user.h and sample implementations can be found in \verb|/example/native.s| and \verb|/example/support.c|
It may also be useful to specify:
\begin{itemize}
\item The object layout for heap objects. See section \ref{sect:objects}
\end{itemize}

Once these are defined, the VM is built as follows, see figure \ref{fig:build}: The \verb|.gsc| files for the interpreter and other support code is created by running \gvmtic{} and \gvmtc{} respectively. These  \verb|.gsc| files are converted to \verb|.gso| with \gvmtas. All the  \verb|.gso| files are then linked together using \glink{} to form as single object file, which can be linked with any native object files to form an executable using the standard linker.

\section{Debugging}
Although the toolkit maintains symbolic information from the source to the executable, some variables may get renamed. 
The signatures of functions may appear different in the debugging, although their names should be unaltered.
The execution stack can be accessed via the stack pointer, the name of which is \verb|gvmt_sp|.

\section{Adding a compiler}
It is probably best to have a working interpreter before adding the compiler. To generate the compiler simply run \gvmtcc{} on the output file from \gvmtic. The resulting C++ file should be compiled and linked using the system C++ compiler. See section \ref{sect:gvmtcc} for more details.

\chapter{User interface}

To create a complete VM, the various components need to interact and to do so correctly.

So this can happen the VM developer is required to implement a number of functions to assist the GVMT. 
The GVMT provides an API for interacting with generated and library components.

Examining the files in the /example subdirectory may help to make things clearer.

\section{User provided code}

The file /include/gvmt/user.h lists all the functions which the developer needs to supply, with the exception of the \verb|gvmt_shape| function which must be compiled natively, not by \gvmtc.

\subsection{Garbage collection\label{sect:user-shape}}
In order for the tooolkit to provide precise garbage collection, the layout of all heap objects must be known. The developer must provide two functions and one integer value to allow GVMT to perform \emph{precise} garbage collection.
\begin{description}
\item [int GVMT\_MAX\_SHAPE\_SIZE] The maximum number of entries required to represent any shape (including terminating zero).
\item [long* gvmt\_shape(GVMT\_Object object, int* buf)] This function returns a vector of integers, representing the shape of the object, as defined in section \ref{sect:shape}. The int array buf can be used to store the shape if required. The buffer, \verb|buf| is guaranteed to include at least GVMT\_MAX\_SHAPE\_SIZE entries.
\item [long gvmt\_length(GVMT\_Object obj)] This function returns the length of the object in bytes.
\item [void user\_finalize\_object(GVMT\_Object o)] This function may be called by the garbage collector for any object that is declared as requiring finalization and is unreachable. Exceptions raised in this function, or any callee, terminate finalisation, but do not propogate. The function is called by the finalizer thread and will be called asynchronously. This function may be called on any or all object declared as finalizable, in any order. 
\end{description}


\subsection{The Marshalling Interface}
In order to support marshalling the user must provide a function \\
\verb|gvmt_write_func get_marshaller_for_object(GVMT_Object object)| to give a write function for each object. The toolkit can create a write function for any object declaration, but it is unable to determine which function should be applied to which object at runtime. 

To assist marshalling to GSC code, the user should also provide the \\
\verb|void gvmt_readable_name(GVMT_Object object, char *buffer)| function to provide meaningful names for objects.
This function should store an ASCII string into the provided buffer. The buffer is guaranteed to be of at least GVMT\_MAX\_OBJECT\_NAME\_LENGTH bytes in length 

\subsection{Debugging}
The GVMT can insert calls to a trace function at the start of each bytecode. Since GVMT knows nothing about the semantics of the VM, the developer must provide the actual trace function. A very simple version is provided in the example.

\section{GVMT provided functions}
GVMT provides a wide range of functions to interact with the underlying abstract machine, as well as a number of intrinsic functions for code written in C.

\subsection{The data stack}
The data stack can be examined and modified in C code, using intinsics:
\begin{itemize}
\item \verb|gvmt_stack_top()| returns a pointer to the top-of-stack.
\item \verb|gvmt_insert(size_t n)| inserts n NULLs onto the stack and returns a pointer to them.
\item \verb|GVMT_PUSH(x)| can be used to push individual values on to the stack.
\item \verb|gvmt_drop(size_t n)| discards n items from the top of the stack.
\end{itemize}
Care should be taken with these functions as they manipulate the data stack, which is used to evaluate expression. It is best to keep any C statement using these intrinsics as simple as possible.

\subsection{The call stack}
The call stack is generally opaque, except for interpreter frames, which can be accessed using the following intinsics:
\begin{itemize}
\item \verb|gvmt_current_frame()| returns a opaque pointer to the currently executing interpreter frame.
\item \verb|gvmt_caller_frame(void* frame)| takes a pointer to an interpreter frame and returns a pointer to its caller.
\item \verb|gvmt_frame_index(char* name)| returns an index for a named local. This index can be used to read the local variables of a given frame.
\item \verb|gvmt_frame_item_r(void* frame, unsigned index)| returns the reference local variable corresponding to the given index.
\item \verb|gvmt_frame_item_p(void* frame, unsigned index)| returns the pointer local variable corresponding to the given index.
\item \verb|gvmt_frame_item_i(void* frame, unsigned index)| returns the integer local variable corresponding to the given index.
\end{itemize}

\subsection{The garbage collector\label{sect:user-roots}}

The following intrinsic allocates object (the GSC code is GC\_MALLOC).
\begin{itemize}
\item \verb|GVMT_Object gvmt_malloc(size_t s)| Allocates a new object. This object should be initialised immediately, and \emph{must} be initialised before the next GC-SAFE point.
\end{itemize}

The garbage collector can automatically find all objects from both stacks and all global variables. Sometimes it is necessary to have some control over the garbage collector.
GVMT provides the following functions to interact with the garbage collector.
\begin{itemize}
\item \verb|gvmt_gc_finalizable(GVMT_Object obj)| Declares that obj will need finalization.
\item \verb|void* gvmt_gc_weak_reference(void)| Returns a weak reference to obj.
\item \verb|GVMT_Object gvmt_gc_read_weak_reference(void* w)| Reads a weak reference, the reference returned will either be the last value written to this weak-reference or NULL.
\item \verb|gvmt_gc_read_weak_reference(void* w, GVMT_Object o)| Writes to a weak reference
\item \verb|void gvmt_gc_free_weak_reference(void* w)| Frees a weak reference if no longer required.
\item \verb|void* gvmt_gc_add_root(void)| Returns a new root to the heap, initialised to obj.
\item \verb|GVMT_Object gvmt_gc_read_root(void* root)| Returns the value referred to by the root.
\item \verb|void gvmt_gc_write_root(void* root, GVMT_Object obj)| Writes a new object to a root.
\item \verb|void gvmt_gc_free_root(void* root)| Deletes this root, the object referred to may now be garbage collected.

\item \verb|void* gvmt_gc_pin(GVMT_Object obj)| Pins the object referred to by \verb|obj|. The garbage collector will not move it. Be aware that collection and pinning are independent. The grabge collector may collect pinned objects. To pass an object to native code, or to use an internal pointer will require both pinning the object \emph{and} retaining a refeerence to it.
\item \verb|gvmt_gc_unpin(GVMT_Object obj)| Unpins the object, which may now by moved by the garbage collector. It is not necessary to unpin objects in order for them to be garbage collected.

\end{itemize}

\subsection{Exception handling\label{sect:user-except}}
Three intrinsics are provided:
\begin{itemize}
\item \verb|GVMT_Object gvmt_protect(void)| Intrinsic for the PROTECT instruction.
\item \verb|void gvmt_unprotect(void)| Intrinsic for the UNPROTECT instruction.
\item \verb|void gvmt_raise(GVMT_Object ex)| Intrinsic for the RAISE instruction.
\end{itemize}

Rather than use \verb|GVMT_Object gvmt_protect()| and \verb|void gvmt_unprotect()| directly developers should use the macros: 
\verb|GVMT_TRY|, \verb|GVMT_CATCH| and \verb|GVMT_END_TRY| as follows:
\begin{verbatim}
GVMT_TRY(exception_variable)
    code which may raise exception
GVMT_CATCH
    code to execute if an exception is raised
GVMT_END_TRY
\end{verbatim}

\subsection{Threading\label{sect:user-threads}}
Since the GVMT relies on the operating system to provide threading support, it must be informed when a new thread is created and used.
To start running GVMT code in a new thread, call:
\begin{verbatim}
int gvmt_enter(uintptr_t stack_space, gvmt_func_ptr func, int pcount, ...);
\end{verbatim}
Where:
\begin{itemize}
\item \verb|stack_space| is the amount of items required on the data stack,
\item \verb|gvmt_func_ptr| is the GSC function to execute, and
\item \verb|pcount| is the number of parameters following.
\end{itemize}

If reentering GVMT from native code that was is running in an already initialised thread, then use 
\begin{verbatim}
int gvmt_reenter(gvmt_func_ptr func, int pcount, ...);
\end{verbatim}
to avoid reinitialising the running thread.

When a thread has finished call \verb|gvmt_finished(void)|.

\appendix

\chapter{Installing the GVMT}

Currently the GVMT has only been tested on the x86-linux platform, but it should work on any posix-compliant x86 platform.
GVMT can be installed in the following steps:
\begin{enumerate}
\item Download GVMT from \url{http://code.google.com/p/gvmt/}
\item Unpack into a clean directory, \$(GVMT)
\item Install llvm (version2.5) from \url{http://llvm.org/releases/download.html\#2.5}
\item Type \verb|make|
\begin{enumerate}
\item[] The build process will fetch lcc using wget. If you do not have wget installed or wget fails then you will have to download lcc from  \url{http://sites.google.com/site/lccretargetablecompiler/downloads/4.2.tar.gz} and save it as \$(GVMT)/lcc.tar.gz
\end{enumerate}
\item Type \verb|sudo make install| (You will need write privileges to /usr/local in order to install the tools)
\end{enumerate}

If it doesn't work, then email me: marks@dcs.gla.ac.uk



\chapter{GSC Instruction set\label{app:inst}}
\small

\input{./instructions.tex}

\end{document}
